\documentclass[11pt,a4paper]{article}
\usepackage[utf8]{inputenc}
\usepackage{amsmath,amssymb,amsthm,mathtools,hyperref,xcolor,geometry}
\geometry{margin=2.5cm}

\newtheorem{theorem}{Theorem}
\newtheorem{proposition}[theorem]{Proposition}
\newtheorem{conjecture}[theorem]{Conjecture}
\theoremstyle{definition}
\newtheorem{remark}[theorem]{Remark}

\newcommand{\AA}{\mathcal{A}}
\newcommand{\R}{\mathbb{R}}
\newcommand{\BI}{\mathrm{B\text{-}I}}
\newcommand{\BIX}{\mathrm{B\text{-}IX}}
\newcommand{\BB}{\mathrm{BB}}

\title{B1 + CCM-BIX revisit: max-effort theoretical re-examination
  in the light of T2-Bianchi I}
\author{Kevin Remondi\`ere\\\small (max-effort attempt by Opus 4.7, 1M ctx)}
\date{3 May 2026}

\begin{document}
\maketitle

\begin{abstract}
We revisit two directions previously marked CLOSED in the algebraic
arrow-of-time programme \cite{Remondiere2026FRW,RemondiereOpus2026Bianchi}:
(i) the FRW apparent-horizon entropy via T2 past asymmetry,
(ii) the Connes--Consani--Moscovici (CCM) zeta-spectral-triple bridge to
Hartnoll--Yang (HY) automorphic dynamics. Today's positive result on
T2-Bianchi I (Theorem~\ref{thm:T2-BI} below, after
\cite{T2BIextension2026}) prompts a re-examination. We find: (B1) NEW
coincidence (algebraic non-cyclic-separating $\Leftrightarrow$ vanishing
horizon entropy) but no new theorem; (CCM-BIX) a sympy-verified
algebraic identification of the CCM dilation with a compactified
Hartnoll--Yang dilation, which re-opens the bridge that yesterday's
\textsc{frw}-grounded NO-GO had closed. Both directions remain partial;
we recommend pursuing CCM-BIX.
\end{abstract}

\section{Inputs}

\begin{enumerate}
  \item \textbf{T2-Bianchi I (today's positive)}
  \cite{T2BIextension2026}. The inductive limit
  $\AA(D_\BB)_\BI = \overline{\bigcup_{t_i\downarrow 0}
   \AA(D_{(t_i,t_f)})_\BI}$ does not admit any Hadamard quasifree state
  (e.g.\ the SLE \cite{BanerjeeNiedermaier2023}) as a cyclic-separating
  vector, in any GNS representation; proof by two independent IR
  divergences (volume-element zero-mode~$\log^3 \delta$ and contracting
  Kasner-direction~$1/|k_1|$).
  \item \textbf{Yesterday's B1 NO-GO}
  \cite{B1FRWentropy2026}. The CLPW
  \cite{CLPW2022} Schwarzschild--de Sitter linearisation step (off-shell
  point energy~$\Rightarrow$ horizon-area shift) has no FRW analog; the
  conformal pullback flattens~$H(t)$ to a universal M\"obius factor~$2\pi$.
  \item \textbf{Yesterday's CCM-FRW NO-GO}
  \cite{CCMFRWbridge2026}. CCM's dilation
  $D = -i u \partial_u$ on $L^2([\lambda^{-1},\lambda], du/u)$
  \cite{CCM2025} is linear in $u$; the FRW dS modular generator
  $K_{\mathrm{FRW}} \sim (R^2 - \eta^2)\partial_\eta$ is quadratic
  in $\eta$. Hence no unitary identification.
\end{enumerate}

\section{B1-revisit}

\subsection{Claim and sympy verification}

Let $(M, g)$ be vacuum Bianchi I (Kasner, $(p_1, p_2, p_3) = (-1/3,
2/3, 2/3)$). Define the cosmological apparent horizon as
$r_{\mathrm{AH}} = 1/H_{\mathrm{avg}}$ with
$H_{\mathrm{avg}} = (1/3) \sum_i (\dot a_i / a_i) = 1/(3t)$. Then:
\begin{equation}
  r_{\mathrm{AH}}(t) = 3t,
  \quad
  A_{\mathrm{AH}}(t) = 4\pi r_{\mathrm{AH}}^2 = 36\pi t^2,
  \quad
  S_{\mathrm{BH}}(t) = \frac{A_{\mathrm{AH}}}{4 G_N} = \frac{9\pi t^2}{G_N}
  \xrightarrow[t\to 0^+]{} 0.
\end{equation}
By Theorem~\ref{thm:T2-BI} (recalled below), the algebra
$\AA(D_\BB)_\BI$ admits no cyclic-separating vector.

\begin{theorem}[T2-Bianchi I; \cite{T2BIextension2026}]\label{thm:T2-BI}
The inductive-limit local algebra
$\AA(D_\BB)_\BI \subset \mathcal{B}(\mathcal{H}_{\mathrm{SLE}})$ does
NOT admit the Banerjee--Niedermaier SLE state, nor any state in its
BFV folium, as a cyclic-separating vector in any GNS representation.
\end{theorem}

\begin{conjecture}[B1-revisit, posited by parent agent]\label{conj:B1-revisit}
The vanishing $S_{\mathrm{BH}}(t) \to 0$ as $t\to 0^+$ AND the algebraic
non-cyclic-separating obstruction of Theorem~\ref{thm:T2-BI} are TWO
MANIFESTATIONS of the same algebraic property of the cosmological
wedge inclusion $\AA(D_\infty) \supset \AA(D_\BB)$.
\end{conjecture}

\subsection{Verdict: PARTIAL}

We sympy-verify (\texttt{B1\_CCM\_BIX\_revisit.py}, Section~1) the
chain:
\begin{itemize}
  \item \emph{Algebraic side.} The smeared Wightman two-point on
  test-functions extending to $t = 0$ contains the zero-mode
  contribution $\bigl(\int_\delta^{2\delta} \log(t)/t\,dt\bigr)^2 \to
  +\infty$, plus the contracting-Kasner-direction $1/|k_1|$ IR.
  Hence no Hadamard normal extension to $\AA(D_\BB)_\BI$.
  \item \emph{Geometric side.} $A_{\mathrm{AH}} \sim t^2$,
  $S_{\mathrm{BH}} \sim t^2/G_N$, both $\to 0$.
  \item \emph{Bekenstein bound (semi-classical).} $\langle K
  \rangle_\omega \le 2\pi r_{\mathrm{AH}} E_{\mathrm{inside}} \le
  S_{\mathrm{BH}}$. Hence
  $\langle K \rangle_{\max} \le S_{\mathrm{BH}}/(2\pi r_{\mathrm{AH}}) =
  3t/(2 G_N) \to 0$.
\end{itemize}

The vanishings COINCIDE at $t \to 0^+$. However:
\begin{itemize}
  \item The Bekenstein bound is \emph{semi-classical}; its derivation
  assumes smooth-spacetime geometry on the diamond, which fails as
  $t\to 0$. So the would-be ``$\langle K\rangle = 0
  \Rightarrow$ no cyclic-separating vector'' is a HEURISTIC, not a
  rigorous deduction. (This is the same caveat as strategy S5 in the
  T2 paper; see \cite{T2BIextension2026} \S~S5.)
  \item There is no operator-level identification of the modular
  Hamiltonian $K$ with an ``area observable'', because the FRW
  pullback flattens the Hubble scale (yesterday's B1 NO-GO,
  \cite{B1FRWentropy2026}~\S~5(a)). T2-Bianchi I uses a direct
  mode-decomposition argument, NOT a pullback, so this objection does
  not kill T2-BI itself; but it does kill the B1-revisit conjecture's
  attempt to use T2-BI as a route to $S = A/4G$.
\end{itemize}

\paragraph{Verdict.} Conjecture~\ref{conj:B1-revisit} is REFUTED as
written: both vanishings have a common geometric cause (the singular
limit), but they are not derivable one from the other within the
present framework. The B1 NO-GO of \cite{B1FRWentropy2026} stands.
\textbf{Re-opens / still closed / partial: PARTIAL} (structural
coincidence noted, no new theorem).

\section{CCM-BIX bridge}

\subsection{Setup}

CCM \cite{CCM2025} eq.~(5.14): on $L^2([\lambda^{-1}, \lambda],
du/u)$ with periodic boundary conditions,
\begin{equation}
  D^{(\lambda)}_{\log} = -i u \frac{\partial}{\partial u},
  \quad
  \mathrm{Spec}\bigl(D^{(\lambda)}_{\log}\bigr) = \biggl\{\frac{2\pi n}{2 \log \lambda} :
   n \in \mathbb Z\biggr\}.
\end{equation}
Theorem~1.1 of \cite{CCM2025}: the rank-one perturbation
$D^{(\lambda, N)}_{\log} = D^{(\lambda)}_{\log} -
|D^{(\lambda)}_{\log} \xi\rangle\langle \delta_N|$
has spectrum (numerically) approaching $\{\gamma_n\}$ as
$N, \lambda \to \infty$.

HY \cite{HartnollYang2025} eq.~(37): in the noncompact picture of the
$SL(2,\R)$ principal series $\Delta = \tfrac{1}{2} + i\varepsilon$,
on $L^2(\R, dx)$,
\begin{equation}
  D \psi = i\Bigl(x \frac{d \psi}{dx} + \Delta \psi\Bigr).
\end{equation}
HY eq.~(48--49): the wavefunction in the dilatation basis is
$\phi(t) = \Phi(\tfrac{1}{2} + it)$ with $\Phi$ proportional to the
$L$-function $\widetilde L(s)$. Vanishing of $\phi(t)$ at the
nontrivial zeros: ``absorption spectrum''.

\subsection{Sympy-verified mapping}

\begin{proposition}[Sympy-verified, \texttt{B1\_CCM\_BIX\_revisit.py}
Section~2]
With $u, x$ identified via the unitary
$U: L^2([\lambda^{-1}, \lambda], du/u) \to L^2([-\log \lambda, \log
\lambda], dx)$ given by $u = e^x$, $du/u = dx$,
\begin{equation}\label{eq:CCM-HY-mapping}
  U D^{(\lambda)}_{\log} U^{-1} = -i \partial_x,
  \qquad
  D_{\mathrm{HY}} = -\,U D^{(\lambda)}_{\log} U^{-1} \cdot x
   + i\Delta \cdot \mathrm{id}.
\end{equation}
The first equation says: CCM's dilation on the
$\log$-rescaled compact interval is the
\textbf{periodic-momentum operator} on $L^2([-\log \lambda, \log \lambda], dx)$.
The second equation, in operator-substituted form, reads
\begin{equation}
  D_{\mathrm{HY}} - i\Delta\,\mathrm{id} = i x \partial_x = -D_{\mathrm{CCM}}
  \;\;(\text{after } u \leftrightarrow x).
\end{equation}
\end{proposition}

\textsc{Sympy verification} (\texttt{B1\_CCM\_BIX\_revisit.py},
output): the difference $D_{\mathrm{HY}}\psi -
i\Delta\psi - (-D_{\mathrm{CCM}})\psi$ is identically zero.

\subsection{Why yesterday's NO-GO does not apply}

Yesterday's NO-GO \cite{CCMFRWbridge2026} compared CCM with the FRW
modular generator $K_{\mathrm{FRW}} = (\pi/2R)(R^2 - \eta^2)
\partial_\eta$, which is quadratic in $\eta$. The HY dilation is
linear (table in \texttt{B1\_CCM\_BIX\_revisit.py} Section~3). Hence:

\begin{itemize}
  \item \textbf{Geometry mismatch.} HY works on Bianchi IX
  (Mixmaster), where the natural dilation is on the BKL Iwasawa
  coordinate (eq.~(1) of \cite{HartnollYang2025}), NOT on FRW
  conformal time. The yesterday NO-GO's quadratic-vs-linear obstruction
  is a verdict on FRW, NOT BIX.
  \item \textbf{Operator match.} On the BIX side, both CCM and HY
  dilations are LINEAR in their respective variables, with
  $D_{\mathrm{HY}} = -D_{\mathrm{CCM}} + i \Delta$ (as
  per~\eqref{eq:CCM-HY-mapping}).
\end{itemize}

\subsection{Residual gaps and verdict}

\begin{itemize}
  \item HY's $L$-function is a wavefunction \emph{overlap}, not an
  eigenvalue. The Riemann zeros are zeros of $\phi(t)$, NOT of
  $D_{\mathrm{HY}}$. (HY \S~6 explicitly: ``zeros = absorption spectrum''.)
  \item CCM's spectrum approaches $\{\gamma_n\}$ only after the
  rank-one perturbation tied to the Euler product over $p \le \lambda^2$.
  Without the perturbation, the spectrum is the trivial torus.
  \item There is no known dictionary mapping CCM's perturbation
  operator to HY's modular-invariance constraint.
  \item HY uses Wheeler--DeWitt on a 1-particle minisuperspace; T2-BI
  uses local-AQFT on a continuum field theory. The frameworks are
  mathematically distinct.
\end{itemize}

\textbf{Verdict CCM-BIX: PARTIAL OPENING.} The dilation-operator
identification is REAL and sympy-verified. But the Hilbert--P\'olya
bridge at the spectral level remains conjectural even after this
identification.

\section{Recommendation}

\begin{enumerate}
  \item \textbf{Pursue CCM-BIX further.} Two specific directions:
  \begin{itemize}
    \item Iwasawa-coordinate dilation on BIX minisuperspace
    (\cite{HartnollYang2025} eq.~(1), Iwasawa $\beta_1, \beta_2,
    \beta_3$): check that the dilation is literally
    $-i \partial_y$ on the upper half-plane, identifying with
    CCM's $-i u \partial_u$ via $u = e^y$. ~1 week.
    \item CCM perturbation $\leftrightarrow$ HY modular constraint
    dictionary: speculative but high-reward. ~2--3 months.
  \end{itemize}
  \item \textbf{Do NOT pursue B1-revisit further.} The coincidence is
  structurally interesting but does not yield a new theorem. The
  original B1 NO-GO stands.
\end{enumerate}

\begin{thebibliography}{99}
\bibitem{Remondiere2026FRW}
K.\ Remondi\`ere, \emph{Algebraic Weyl Curvature Hypothesis: FRW Theorem
T1+T2+T3}, companion note \texttt{algebraic\_arrow.tex} (May 2026).
\bibitem{RemondiereOpus2026Bianchi}
K.\ Remondi\`ere (with max-effort attempt by Opus 4.7), \emph{Bianchi
extension attempt: Proposition B.1 (T1 perturbative lift)}, companion note
\texttt{bianchi\_extension\_opus.tex} (May 2026).
\bibitem{T2BIextension2026}
K.\ Remondi\`ere (with Opus 4.7), \emph{T2 Bianchi extension},
\texttt{/tmp/T2\_bianchi\_extension.\{tex,md,py\}} (May 2026).
\bibitem{B1FRWentropy2026}
K.\ Remondi\`ere (with Opus 4.7), \emph{B1 -- Witten/CLPW
generalised-entropy formula on FRW type II$_\infty$ algebras},
\texttt{/tmp/B1\_frw\_entropy.md} (May 2026).
\bibitem{CCMFRWbridge2026}
K.\ Remondi\`ere (with Opus 4.7), \emph{CCM 2511.22755 + Hartnoll--Yang
2502.02661 vs FRW Type II$_\infty$},
\texttt{/tmp/ccm\_hartnoll\_frw\_bridge.md} (May 2026).
\bibitem{CLPW2022}
V.\ Chandrasekaran, R.\ Longo, G.\ Penington, E.\ Witten, \emph{An algebra
of observables for de Sitter space}, JHEP 02 (2023) 082, arXiv:2206.10780.
\bibitem{CCM2025}
A.\ Connes, C.\ Consani, H.\ Moscovici, \emph{Zeta Spectral Triples},
arXiv:2511.22755 (2025).
\bibitem{HartnollYang2025}
S.A.\ Hartnoll, M.\ Yang, \emph{The Conformal Primon Gas at the End
of Time}, arXiv:2502.02661 (2025).
\bibitem{BanerjeeNiedermaier2023}
R.\ Banerjee, M.\ Niedermaier, \emph{States of Low Energy on Bianchi I
spacetimes}, J.\ Math.\ Phys.\ \textbf{64}, 113503 (2023),
arXiv:2305.11388.
\end{thebibliography}

\end{document}
